%! AP Statistics — Unit 3, Lesson 3.1 — Lesson Plan
\documentclass[10pt]{article}
\usepackage{apstats-article}
\usepackage{apstats-boxes}

\newcommand{\UnitNumberName}{Unit 3: Collecting Data \quad}
\newcommand{\LessonNumberName}{Lesson 3.1: Introducing Statistics --- Do the Data We Collected Tell the Truth?}

\begin{document}
\begin{center}
    {\Large\bfseries \CourseName: \SchoolYear\ (\MeetingLength) \\
    {\small\bfseries \UnitNumberName \LessonNumberName}}
\end{center}

\begin{tcolorbox}[colback=sky, colframe=navy, boxrule=0.9pt, arc=2mm,
    left=2mm, right=2mm, top=1.5mm, bottom=1.5mm]
    \textbf{Primary Objective:} Students will recognize that the method used to collect
    data determines the trustworthiness of conclusions --- methods that do not rely on
    chance can produce biased, untrustworthy results (Big Idea: VAR).
\end{tcolorbox}

\renewcommand{\arraystretch}{1.6}
\begin{skillbox}[Priority Ideas \& Skills]{goldbox}
\begin{minipage}[t]{0.48\linewidth}
    \textbf{AP Skill 1 --- Selecting Statistical Methods}
    \begin{itemize}
        \item Identify a statistical question about a data collection method (Skill 1.A).
        \item Recognize that methods not relying on chance produce untrustworthy
              conclusions (VAR-1.E, VAR-1.E.1).
        \item Transition from Unit 2 (describing relationships) to Unit 3 (how data
              are gathered).
    \end{itemize}
\end{minipage}
\hfill
\begin{minipage}[t]{0.48\linewidth}
    \textbf{Key Understandings}
    \begin{itemize}
        \item The way data are collected determines what conclusions are valid.
        \item ``Convenient'' data often reflect the preferences of who chose to respond,
              not the broader population.
        \item Chance-based methods are the foundation of trustworthy statistics.
        \item This lesson sets up the entire unit: sampling (3.2--3.4) and
              experiments (3.5--3.7).
    \end{itemize}
\end{minipage}
\end{skillbox}

\begin{skillbox}[{Vocabulary, Concepts, \& Theorems}]{greenbox}
\begin{minipage}[t]{0.48\linewidth}
    \begin{itemize}
        \item \textbf{Data collection method} --- the procedure used to gather information
        \item \textbf{Chance-based method} --- selection that relies on a random mechanism
        \item \textbf{Non-chance method} --- selection based on convenience, judgment, or
              voluntary participation
    \end{itemize}
\end{minipage}
\hfill
\begin{minipage}[t]{0.48\linewidth}
    \begin{itemize}
        \item \textbf{Trustworthy conclusion} --- a result that can be generalized to the
              population of interest
        \item \textbf{Bias} --- a systematic tendency for the data to misrepresent the
              population
        \item \textbf{Anecdotal evidence} --- informal evidence based on a few cases, not
              a systematic sample
    \end{itemize}
\end{minipage}
\end{skillbox}

\begin{skillbox}[Activate Prior Knowledge \& Spiral Review]{sky}
\begin{minipage}[t]{0.48\linewidth}
    \textbf{Prior Knowledge Check (3--4 min)}
    \begin{itemize}
        \item Quick review: ``What is a confounding variable?'' (Unit 2)
        \item Connect: Unit 2 found associations --- Unit 3 asks how we should collect
              data so our conclusions mean something.
        \item Prompt: \textit{``If I want to know whether energy drinks improve test
              scores, how should I collect the data?''}
    \end{itemize}
\end{minipage}
\hfill
\begin{minipage}[t]{0.48\linewidth}
    \textbf{Connections Forward}
    \begin{itemize}
        \item Observational studies vs.\ experiments $\to$ Lesson 3.2
        \item Random sampling methods $\to$ Lesson 3.3
        \item Bias in non-random samples $\to$ Lesson 3.4
        \item Experimental design $\to$ Lessons 3.5--3.6
        \item Causal inference $\to$ Lesson 3.7
    \end{itemize}
\end{minipage}
\end{skillbox}

\begin{skillbox}[Hook \& Lesson]{sky}
\begin{minipage}[t]{0.48\linewidth}
    \textbf{Hook (5--7 min)}\\
    Display three ``survey results'':
    \begin{itemize}
        \item A grocery store places a comment card near the exit --- 12 customers
              return it, all complaining about prices.
        \item A news website polls online visitors about a new school bond.
        \item A researcher randomly dials 400 phone numbers.
    \end{itemize}
    Ask: ``Which result would you trust most? Why?''\\
    Transition: \textit{``The answer depends entirely on HOW the data were collected.''}
\end{minipage}
\hfill
\begin{minipage}[t]{0.48\linewidth}
    \textbf{Lesson Flow (15--18 min)}
    \begin{enumerate}
        \item Define data collection method; show that the same question yields
              different answers depending on who is asked and how.
        \item Contrast: ``I asked my friends'' vs.\ ``I used a random number generator.''
        \item Key idea (VAR-1.E.1): non-chance methods produce results that may
              systematically misrepresent the population.
        \item Counter-intuitive point: a large non-random sample is \emph{not} more
              trustworthy than a small random one (Literary Digest 1936 example).
        \item Preview road map: rest of Unit 3 develops chance-based and well-designed
              methods for both surveys and experiments.
    \end{enumerate}
\end{minipage}
\end{skillbox}

\begin{skillbox}[Explicit Instruction \& Active Monitoring]{sky}
\begin{minipage}[t]{0.48\linewidth}
    \textbf{Direct Instruction (10 min)}
    \begin{itemize}
        \item Work through 3 data collection scenarios. For each: (a) Was chance used?
              (b) Whose opinions are represented? (c) Can we trust the conclusion for
              the whole population?
        \item Emphasize VAR-1.E.1: non-chance methods produce untrustworthy conclusions
              because the sample may not represent the population.
    \end{itemize}
\end{minipage}
\hfill
\begin{minipage}[t]{0.48\linewidth}
    \textbf{Active Monitoring}
    \begin{itemize}
        \item Listen for: ``bigger sample = more trustworthy'' --- address with the
              Literary Digest example.
        \item Cold-call: ``Tell me one reason the online poll result might not represent
              the population.''
        \item Circulate during guided notes; check students can distinguish chance from
              non-chance selection.
    \end{itemize}
\end{minipage}
\end{skillbox}

\begin{skillbox}[Group Work \& Differentiation]{redbox}
\begin{minipage}[t]{0.48\linewidth}
    \textbf{Partner Activity (8--10 min)}\\
    Four data collection scenarios. Partners decide:
    \begin{enumerate}
        \item Was chance used to select individuals?
        \item Who is represented and who is left out?
        \item Can conclusions be trusted for the population?
        \item What one change would make the data more trustworthy?
    \end{enumerate}
\end{minipage}
\hfill
\begin{minipage}[t]{0.48\linewidth}
    \textbf{Differentiation}
    \begin{itemize}
        \item \textit{Support}: Anchor chart --- ``Was chance used? YES $\to$ more
              trustworthy. NO $\to$ less trustworthy.''
        \item \textit{Extension}: Research the 1936 Literary Digest poll. Write a
              paragraph explaining why 2.4 million surveys produced the wrong prediction.
        \item \textit{ELL}: Vocabulary card with ``chance,'' ``sample,'' and
              ``population'' in accessible language.
    \end{itemize}
\end{minipage}
\end{skillbox}

\begin{skillbox}[Individual Work \& Assessment]{redbox}
\begin{minipage}[t]{0.48\linewidth}
    \textbf{Exit Ticket (5 min)}\\
    Scenario: \textit{``A principal wants to know how students feel about the new lunch
    menu. She asks all 30 students in her AP Statistics class to fill out a survey.''}
    \begin{enumerate}
        \item Was chance used to select these students?
        \item Why might these results not represent all students?
        \item Suggest one change that would make the data more trustworthy.
    \end{enumerate}
\end{minipage}
\hfill
\begin{minipage}[t]{0.48\linewidth}
    \textbf{AP-Style Check}\\
    \textit{``A TV station asks viewers to call in and vote on a new city ordinance.
    Of 3{,}000 callers, 78\% oppose it. Which best explains why this may not be a
    trustworthy estimate of public opinion?''}
    \begin{enumerate}[label=(\Alph*)]
        \item The sample size is too small.
        \item Only viewers with strong opinions are likely to call in.
        \item The ordinance is too controversial to survey fairly.
        \item Television viewers are a representative cross-section of the public.
    \end{enumerate}
    Answer: (B)
\end{minipage}
\end{skillbox}

\begin{skillbox}[Reinforcement \& Extension]{goldbox}
\begin{minipage}[t]{0.48\linewidth}
    \textbf{Homework / Practice}
    \begin{itemize}
        \item For 4 real-world scenarios, identify whether chance was used, who is
              represented, and whether results are trustworthy.
        \item Reflection: \textit{``Why isn't a large sample size enough to guarantee
              trustworthy data?''}
    \end{itemize}
\end{minipage}
\hfill
\begin{minipage}[t]{0.48\linewidth}
    \textbf{Extension / Preview}
    \begin{itemize}
        \item \textit{Extension}: Look up the 1936 Literary Digest poll. Explain why
              10 million surveys led to the wrong prediction.
        \item \textit{Preview}: Lesson 3.2 introduces the formal distinction between
              observational studies and experiments --- start thinking about what
              ``imposing a treatment'' means.
    \end{itemize}
\end{minipage}
\end{skillbox}

\end{document}
